\section{The Propers: Themes and Movement}
\zlabel{triptych:concise:themes:start}
Christ's word makes faithful action possible before the outcome is seen;
God's promise makes thanksgiving possible while suffering remains. The
father who starts home and the captives who remember Zion show two forms
of this hope. One has received assurance of a particular healing. The
others still weep. The Sunday gives both a place within the same worship.

\subsection{Need confessed, service made possible}
The Introit begins with judgment and mercy. Azarias speaks from a furnace
into which fidelity, not idolatry, has brought him. Jerome explains his
confession as the people's voice: the innocent youths acknowledge the
sins of the community to which they belong. Their prayer therefore
joins solidarity to repentance without making each sufferer the author
of his suffering. The psalm verse turns toward the blessed life of
obedience, and the Collect asks for the pardon and peace that enable it.
The assembly seeks a secure mind in order to serve, not an exemption
from responsibility.\footnote{Jerome, \work{Commentary on Daniel}, I,
at 3:29; Schuster, \work{The Sacramentary}, III, p.~175.}

This opening gives both trust and common life their proper beginning.
The petitioner need not pretend to be self-sufficient before asking for
help. The community need not pretend to be innocent of every division
before asking to serve together. In each case mercy comes first. The
Postcommunion will return to the commandments as a gift sought from God:
obedience is the fruit for which the worshippers pray, rather than a
price they announce themselves able to pay.

\subsection{A word received, a household changed}
The Gospel makes dependence concrete. A father asks Jesus to come down
before his son dies. Christ answers by sending the father home, and
the man goes because he believes the word. The servants' report then
confirms the child's recovery at the corresponding hour; belief embraces
the whole household. Trust has become action before the father possesses
that confirmation. The confirmation nevertheless matters: John records
both the inquiry and the answer.

Gregory sees an imperfect believer whose confidence in Christ's healing
still depends on bodily presence. Augustine judges the initial approach
more severely as coldness or unbelief, against the Samaritans' readiness
to hear. They agree that Christ's power exceeds the demand placed on it,
but differ about the father's starting point. The word educates faith
while the healing restores a real child. The household's belief extends
the mercy beyond the solitary petitioner; the servants themselves carry
testimony on which the father acts.\footnote{Gregory, \work{Homilies on
the Gospels}, II.28.1--2; Augustine, \work{Tractates on John}, XVI.1--5.}

Ephesians describes the less dramatic daily shape of such trust: careful
walking, discernment, inward attention and mutual submission. Chrysostom
reads redeeming time as prudent fidelity in hostile circumstances, even
at the cost of relinquishing an advantage. The same Lord who receives
the father's urgent petition teaches a community to act without letting
anger or appetite determine its next step. Singing in the heart is
attentive speech, and the appointed Epistle ends with service to one
another. Faith's inward life has outward consequences.

This movement also changes the meaning of strength. The ruler must
receive a word instead of securing the healer's bodily attendance;
the prosperous must learn to serve rather than expect service alone.
Gregory's contrast with Christ's care for a centurion's servant exposes
the error of honoring possessions more than the person. The healed
household and the mutually serving assembly both receive their good
from a Lord whose regard is not measured by rank. Dependence on him
can free them from dependence on the honors they confer upon one another.

\newpage
\subsection{Food, song and tears on the same journey}
The Gradual praises the giver who opens his hand to all. Augustine's
physician image addresses the person whose request has not received an
immediate answer: care includes knowledge of the fitting time and good.
Bellarmine addresses those who can give: human hoarding contradicts
the generosity praised. These accents belong together. Patience in
need and responsibility toward a hungry neighbor are different demands;
the latter cannot be evaded by quoting the former. Schuster carries
the chant's nourishment toward the Eucharist and finally the sight of
God. Creaturely food remains a good even as desire reaches beyond it.
\footnote{Augustine, Ps.~144, \S\S14--17; Bellarmine, Ps.~144:14--17;
Schuster, III, pp.~175--176.}

The Alleluia's prepared heart sings, while the Offertory's captives sit
and weep. Augustine reads Babylon's passing currents as the loves that
threaten to carry the pilgrim away. Remembering Zion keeps desire fixed
on an abiding home. His exposition ends by inviting song among the
pilgrims themselves. Thus the tears do not cancel Ephesians' command
to sing to one another. They distinguish hope from satisfaction with
the present order. Thanksgiving need not declare captivity pleasant
or bodily pain unreal.\footnote{Augustine, Ps.~136, \S\S2--9,13.}

Such remembrance need not remove the pilgrim from ordinary responsibilities.
Bellarmine distinguishes inward allegiance from outward occupation:
one may hold authority as a burden of service rather than an object
of self-admiration. The open hand and the remembered home then belong
to the same life. Desire for Zion changes how one possesses and gives
in Babylon. It neither identifies every present good with the final
home nor abandons the neighbors who depend on those goods now.

Chrysostom's thanksgiving extends into affliction; the NPNF editor's
note restricts ``all things'' to blessings. Chrysostom's stronger demand
rests on God's goodness, not on approval of wrongdoing. He locates evil
in conduct rather than in created time itself. The Introit confesses
sin, and the Secret asks for vice to be expelled: gratitude cannot
make those same vices praiseworthy. The father seeks help for the sick,
and the assembly brings its lament to the altar.

\subsection{Heavenly medicine, remembered promise, enduring peace}
The Secret asks for healing at the roots of the heart's disorder.
Schuster expounds the mysteries as Eucharistic medicine. Pride, avarice
and resentment injure the petitioner and fracture common life; the
remedy restores the capacity to obey and love. The Communion then
asks God to remember the word that has given hope. Augustine explains
remembrance as fulfilment of the promise, not repair of divine memory.
The promise comforts in humiliation and reaches toward resurrection.
Schuster's liturgical reading identifies the Word with Christ, the
ground of the Church's hope.\footnote{Augustine, Ps.~118, \S\S50--56;
Schuster, III, pp.~177--178.}

When the Maternity prayers are commemorated, the Collect confesses
the Word's real birth from Mary and asks her intercession. Its Secret
joins present to perpetual peace; its Postcommunion asks cleansing and
fellowship in heavenly remedy. These petitions deepen the same dependence
on Christ. The Trinity Preface remains the assembly's thanksgiving in
either branch: Father, Son and Holy Spirit are adored in one essence,
distinct Persons and equal majesty. The journey of trust and the song
of pilgrimage end in God himself. That end gives present gifts their
promise and present service its lasting seriousness.
\zlabel{triptych:concise:themes:end}
\clearpage
